% (c) Johan Hoffman and Anders Logg 2003.
% Licensed under the GNU GPL Version 2.

\documentclass[12pt]{article}

%--- Packages ---
\usepackage[latin1]{inputenc}
\usepackage{graphicx}
\usepackage{fancyhdr}

\setlength{\parindent}{0pt}
\setlength{\parskip}{12pt}
\sloppy

%--- Header an footer ---
\lhead{\sf Puffin}
\rhead{\sf User Manual}
\pagestyle{fancy}
\renewcommand{\footrulewidth}{1pt}
\renewcommand{\headrulewidth}{1pt}

%--- Commands ---
\newcommand{\fixme}[1]{\ \\ \begin{tabular}{|c|}\hline \rm \textbf{FIXME:} \rm  #1  \\ \hline \end{tabular} \\}

\begin{document}

%--- Title page ---

\thispagestyle{empty}

\Large
\textbf{Puffin User Manual} \\
\vspace{-0.5cm}
\hrule
\hfill
\small
\today

\vspace{3cm}

\begin{center}
\includegraphics[width=6cm]{eps/puffin.eps}
\end{center}

\vfill

\normalsize
\textbf{Johan Hoffman and Anders Logg}
\hrule

\newpage

%--- List of authors ---

\ 
\vfill
This manual is written by: \\
Johan Hoffman and Anders Logg.
\thispagestyle{empty}

\cleardoublepage

%--- Table of contents ---

\tableofcontents
\newpage

%--- New chapter ---
\cleardoublepage
\section{Introduction}

This is a first draft for a Puffin manual. Contributions are most welcome.

Puffin is a simple and minimal implementation of FEniCS for Octave (MATLAB).
It is based around the two functions \texttt{AssembleMatrix()} and \texttt{AssembleVector()},
which are used to assemble a linear system
\begin{equation}
  AU = b,
\end{equation}
representing a variational formulation of a differential equation,
\begin{equation}
  a(u,v;w) = l(v;w) \quad \forall v \in V,
\end{equation}
where $a(u,v;w)$ is a bilinear form in $u$ (the trial function) and $v$ (the test function),
and $l(v;w)$ is a linear form in $v$. The bilinear form $a$ and the linear form $l$ are allowed
to depend on a vector $w$ (the coefficients), and are referred to as the left-hand side and
the right-hand side of the variational formulation.

%--- New chapter ---
\cleardoublepage
\section{Mesh format}

The mesh is assumed to be given by the three matrices
\texttt{points}, \texttt{edges}, and \texttt{triangles}, representing the
\emph{points}, \emph{edges}, and \emph{triangles} of an unstructured
triangular mesh in the plane.

\fixme{Add detailed description of the format.}

%--- New chapter ---
\cleardoublepage
\section{Specifying the variational formulation}

A variational formulation is given as a function in the following way:

\begin{verbatim}
function integral = MyForm(u, v, w, du, dv, dw, dx, ds, x, d, t, eq)

if eq == 1
  integral = ... * dx + ... * ds;
else
  integral = ... * dx + ... * ds;
end
\end{verbatim}

The output argument \texttt{integral} should be the integral of the left-hand or right-hand side
(depending on the value of \texttt{eq}) of the variational formulation over a triangle (with 
area \texttt{dx}) or an edge (with length \texttt{ds}).

The input arguments are given by
\begin{itemize}
\item \texttt{u}:  the value of the trial function $u$,
\item \texttt{v}:  the value of the test function $v$,
\item \texttt{w}:  a vector of coefficient values (optional),
\item \texttt{du}: the gradient $\nabla u$ of the trial function $u$,
\item \texttt{dv}: the gradient $\nabla v$ of the test function $v$,
\item \texttt{dw}: a matrix of coefficient gradients (optional),
\item \texttt{dx}: the area of the current triangle (zero if we are on an edge),
\item \texttt{ds}: the length of the current edge (zero if we are on a triangle),
\item \texttt{x}:  the current quadrature point,
\item \texttt{d}:  the current domain number or edge number,
\item \texttt{t}:  the value of time,
\item \texttt{eq}: specifies left-hand side (\texttt{1}) or right-hand side (\texttt{2}) of the variational formulation.
\end{itemize}

As an example, consider the specification of the variational formulation for Poisson's equation.
The variational formulation in mathematical notation is given by
\begin{equation} \label{eq:varform}
  \int_{\Omega} \nabla u \cdot \nabla v \, dx +
  \int_{\Gamma} \gamma u v \, ds =
  \int_{\Omega} fv \, dx + \int_{\Gamma} (\gamma g_D - g_N) v \, ds 
  \quad \forall v.
\end{equation}
For a certain choice of $a$, $f$, $\gamma$, $g_D$, and $g_N$, this can be specified  by a function \texttt{Poisson()} as follows.

\begin{verbatim}
function integral = Poisson(u, v, w, du, dv, dw, dx, ds, x, d, t, eq)

if eq == 1
  integral = du'*dv*dx + g(x,d,t)*u*v*ds;
else
  integral = f(x,d,t)*v*dx + (g(x,d,t)*gd(x,d,t) - gn(x,d,t))*v*ds;
end

%--- Conductivity (penalty factor) ---
function y = g(x, d, t)
y = 1e7;

%--- Dirichlet boundary condition ----
function y = gd(x, d, t)
y = 0;

%--- Neumann boundary condition ---
function y = gn(x, d, t)
y = 0;

%--- Right-hand side, source term ---
function y = f(x, d, t)
y = 5*pi^2*sin(pi*x(1))*sin(2*pi*x(2));
\end{verbatim}

%--- New chapter ---
\cleardoublepage
\section{The function \texttt{AssembleMatrix()}}

The syntax of the function \texttt{AssembleMatrix()} is

\quad \texttt{A = AssembleMatrix(points, edges, triangles, pde, W, time)},

where the output argument is \texttt{A}, the matrix to be assembled, and
the input arguments are given by

\begin{itemize}
\item \texttt{points}:    the matrix containing the node coordinates of the mesh,
\item \texttt{edges}:     the matrix containing the edges of the mesh,
\item \texttt{triangles}: the matrix containing the triangles of the mesh,
\item \texttt{pde}:       the name of the function specifying the variational formulation,
\item \texttt{W}:         a matrix where each column contains the nodal values of additional functions (coefficients) appearing in the
                          variational formulation as the functions \texttt{w(1)}, \texttt{w(2)} and so on (leave empty (\texttt{[]}) if not needed),
\item \texttt{time}:      the value of time (the variable \texttt{t} in the variational formulation).
\end{itemize}

%--- New chapter ---
\cleardoublepage
\section{The function \texttt{AssembleVector()}}

The syntax of the function \texttt{AssembleVector()} is

\quad \texttt{b = AssembleVector(points, edges, triangles, pde, W, time)},

where the output argument is \texttt{b}, the vector to be assembled, and
the input arguments are given by

\begin{itemize}
\item \texttt{points}:    the matrix containing the node coordinates of the mesh,
\item \texttt{edges}:     the matrix containing the edges of the mesh,
\item \texttt{triangles}: the matrix containing the triangles of the mesh,
\item \texttt{pde}:       the name of the function specifying the variational formulation,
\item \texttt{W}:         a matrix where each column contains the nodal values of additional functions (coefficients) appearing in the
                          variational formulation as the functions \texttt{w(1)}, \texttt{w(2)} and so on (leave empty (\texttt{[]}) if not needed),
\item \texttt{time}:      the value of time (the variable \texttt{t} in the variational formulation).
\end{itemize}

%--- New chapter ---
\cleardoublepage
\section{Example programs}

\subsection{The program \texttt{PoissonSolver}}

The program \texttt{PoissonSolver} solves Poisson's equation on the unit square,
using the variational formulation specified by the function \texttt{Poisson()}.

\subsection{The program \texttt{ConvDiffStationarySolver}}

The program \texttt{ConvDiffStationarySolver} solves the stationary convection-diffusion equation around a cylinder,
using the variational formulation specified by the function \texttt{ConvDiffStationary()}.

\subsection{The program \texttt{ConvDiffTimeDepSolver}}

The program \texttt{ConvDiffTimeDepSolver} solves the time-dependent convection-diffusion equation on the unit square,
using the variational formulation specified by the function \texttt{ConvDiffTimeDep()}.

\end{document}
